\documentclass[11pt]{article}
\usepackage[margin=1in]{geometry}
\usepackage{amsmath,amssymb}
\usepackage{bm}
\usepackage{hyperref}
\usepackage{enumitem}

\begin{document}

I’ll go model by model and translate the code into math, plus what it means physically for your signal-loss analysis. I’ll use consistent notation so it’s easier to see how the pieces fit together.

\section*{0. Data and notation (what the arrays actually are)}

From the code:

\begin{itemize}[leftmargin=*]
  \item \texttt{uvp\_power\_list} $\Rightarrow$ list over redundant groups of
  \[
    P_{\mathrm{coh}}(t): \text{ coherent PSPEC at delay } \approx 0, \text{ as a function of LST.}
  \]

  \item \texttt{uvpspec\_avg\_power\_list} $\Rightarrow$ list of
  \[
    P_{\mathrm{inc}}(t): \text{\emph{incoherently} averaged PSPEC at delay } \approx 0, \text{ same LSTs.}
  \]

  \item \texttt{uvpspec\_noise\_list} $\Rightarrow$ list of
  \[
    P_N(t): \text{ noise power estimate from \texttt{uvpspec\_averaged.get\_stats('autos\_diag', key)} at delay } \approx 0.
  \]

  \item \texttt{uvpspec\_signalnoise\_list} $\Rightarrow$ list of
  \[
    P_{SN}(t): \text{ some signal+noise RMS combination you construct from } P_S \text{ and } P_N.
  \]

  \item \texttt{uvpspec\_SNR\_list} $\Rightarrow$ list of
  \[
    \mathrm{SNR}(t) = \frac{\bigl|P_{\mathrm{coh}}(t)\bigr|}{P_N(t)} \quad \text{at delay } \approx 0.
  \]

  \item Each \texttt{list} entry corresponds to one redundant group $g$. Inside each group you have samples indexed by LST $t$.
\end{itemize}

I’ll denote:

\begin{itemize}[leftmargin=*]
  \item Group index: $g = 1, \dots, G$,
  \item Time/LST index within group: $i = 1, \dots, N_g$.
\end{itemize}

So for each group and time you have:

\begin{itemize}[leftmargin=*]
  \item $P_{\mathrm{coh}}^{(g)}(t_i)$,
  \item $P_{\mathrm{inc}}^{(g)}(t_i)$,
  \item $P_N^{(g)}(t_i)$,
  \item $P_{SN}^{(g)}(t_i)$,
  \item $\mathrm{SNR}^{(g)}(t_i)$.
\end{itemize}

The global concatenation across all groups has total length
\[
  N = \sum_g N_g.
\]

\section*{1. ``Bump and valley'' segmentation in LST}

\subsection*{Function: \texttt{bump\_n\_val\_segments\_by\_dips}}

\textbf{Goal:} segment the LST axis into ``bumps'' separated by ``dips'' in the coherent PSPEC, and identify valley regions around each dip.

\textbf{Inputs:}

\begin{itemize}[leftmargin=*]
  \item \texttt{lst}: $t_i$, monotonically increasing LST (in hours, ``rolled'').
  \item \texttt{pspec}: main 1D power spectrum vs.\ LST (e.g.\ coherent power).
  \item \texttt{pspec2}: a second spectrum sampled at the same LST points (e.g.\ incoherent power).
\end{itemize}

\subsubsection*{Step 1 -- smoothing}

\begin{verbatim}
psmooth = savgol_filter(pspec, smooth_window, polyorder)
\end{verbatim}

Mathematically: $psmooth(t)$ is a Savitzky--Golay smoothed version of $pspec(t)$, i.e.\ a local polynomial fit of degree \texttt{polyorder} over a window of length \texttt{smooth\_window}. This reduces small-scale noise so that dips/peaks are better defined.

\subsubsection*{Step 2 -- find dips}

\begin{verbatim}
dip_indices, _ = find_peaks(-psmooth,
                            prominence=dip_prominence,
                            width=dip_width)
\end{verbatim}

\texttt{find\_peaks(-psmooth)} looks for \emph{peaks} in $-psmooth$, i.e.\ \emph{local minima in \texttt{psmooth}}.

Each detected dip index $d_k$ satisfies:

\begin{itemize}[leftmargin=*]
  \item $psmooth$ is locally low compared to its neighbours.
  \item ``Prominence'' and ``width'' ensure the dip is not just noise.
\end{itemize}

\subsubsection*{Step 3 -- segment into ``bumps''}

\begin{verbatim}
bounds = np.concatenate(([0], dip_indices, [len(pspec)-1]))
segments = []
for i in range(len(bounds)-1):
    i0, i1 = bounds[i], bounds[i+1]
    segments.append((lst[i0:i1+1], pspec[i0:i1+1]))
\end{verbatim}

You form boundaries
\[
  0 = b_0 < b_1 = d_1 < b_2 = d_2 < \dots < b_K = d_K < b_{K+1} = N-1.
\]

Each consecutive interval $[b_i, b_{i+1}]$ is a ``bump'' segment:

\[
  \mathcal{S}_i = \{(t_j, pspec_j): j \in [b_i, b_{i+1}]\}.
\]

Physically: each bump is a contiguous LST region between dips -- likely corresponding to a strong sky feature or a beam/sidelobe region.

\subsubsection*{Step 4 -- find peaks and valley ranges around each dip}

\begin{verbatim}
peak_indices, _ = find_peaks(psmooth)
\end{verbatim}

Now you search for \emph{peaks} (local maxima) in $psmooth$. So each \texttt{dip} is flanked by a left peak \texttt{pl} and right peak \texttt{pr}.

For each dip index \texttt{pd}:

\begin{enumerate}[leftmargin=*]
  \item Identify nearest left and right peaks:
\begin{verbatim}
lefts  = peak_indices[peak_indices < pd]
rights = peak_indices[peak_indices > pd]
pl, pr = lefts.max(), rights.min()
\end{verbatim}

So $pl < pd < pr$ with $pl$ and $pr$ the closest peaks.

  \item Define threshold levels as fractions of the peak heights:
\begin{verbatim}
thr_l = frac * pspec[pl]
thr_r = frac * pspec[pr]
\end{verbatim}
So
\[
  T_L = f \cdot pspec(pl), \quad T_R = f \cdot pspec(pr),
\]
with \texttt{frac = f} (e.g.\ 0.05 or 0.25).

  \item Find the valley ``extent'' where the spectrum falls below these thresholds.
  
  On the left side:
\begin{verbatim}
left_ok = np.nonzero(pspec[:pd] > thr_l)[0]
start = (left_ok.max()+1) if left_ok.size else 0
\end{verbatim}
This finds the last index before the dip where $pspec > T_L$. The valley begins at the next index \texttt{start}.

  On the right side:
\begin{verbatim}
right_ok = np.nonzero(pspec[pd+1:] > thr_r)[0]
end = (pd + right_ok.min()) if right_ok.size else len(pspec)-1
\end{verbatim}
This finds the first index after the dip where $pspec > T_R$. The valley ends at this \texttt{end}.

The valley range is
\[
  i \in [\mathrm{start}, \mathrm{end}].
\]

  \item Store the valley in both \texttt{pspec} and \texttt{pspec2}:
\begin{verbatim}
valley_ranges.append((lst[start:end+1], pspec[start:end+1]))
valley_ranges2.append((lst[start:end+1], pspec2[start:end+1]))
\end{verbatim}
\end{enumerate}

\textbf{Physical use:} in your plotting code you \emph{could} use these valley slices to focus on low-power regions between bright ``bumps'' when studying loss vs.\ LST or forming ratios $P_{\mathrm{coh}} / P_{\mathrm{inc}}$.

\section*{2. Two-segment piecewise linear ``slope'' model (signal-loss as a break in log--log space)}

\subsection*{Functions: \texttt{piecewise\_linear}, \texttt{manual\_piecewise\_linear\_fit}, \texttt{joint\_piecewise\_linear\_fit}}

\subsubsection*{Model definition}

\begin{verbatim}
def piecewise_linear(params, x):
    m1, b1, m2, b2, x0 = params
    return np.where(x <= x0,
                    m1 * x + b1,
                    m2 * x + b2)
\end{verbatim}

Given parameters $(m_1, b_1, m_2, b_2, x_0)$, the model is
\[
  f(x) =
  \begin{cases}
    m_1 x + b_1, & x \le x_0, \\
    m_2 x + b_2, & x > x_0.
  \end{cases}
\]

\subsubsection*{Residuals and cost}

For data points $(x_i, y_i)$:

\begin{verbatim}
def residuals(params, x, y):
    return piecewise_linear(params, x) - y
\end{verbatim}

Residual for point $i$ is
\[
  r_i = f(x_i) - y_i.
\]

Sum of squared errors (SSE):
\[
  \mathrm{SSE}(\text{params}) = \sum_i r_i^2.
\]

\subsubsection*{Manual grid search fit}

\texttt{fit\_piecewise\_linear} always forces \texttt{success=False}, so you always use:

\begin{verbatim}
def manual_piecewise_linear_fit(x, y, num_breaks=200):
    candidates = np.linspace(np.min(x), np.max(x), num_breaks)
    best = {'cost': np.inf}
    for x0 in candidates:
        left  = x <= x0
        right = x > x0
        if np.sum(left) < 2 or np.sum(right) < 2:
            continue
        m1, b1 = np.polyfit(x[left], y[left], 1)
        m2, b2 = np.polyfit(x[right], y[right], 1)
        ...
        cost = np.sum((y_pred - y)**2)
        if cost < best['cost']:
            best.update(...)
\end{verbatim}

For each candidate breakpoint $x_0$:

\begin{enumerate}[leftmargin=*]
  \item Split data into left and right sets:
  \[
    \mathcal{L} = \{i : x_i \le x_0\}, \quad
    \mathcal{R} = \{i : x_i > x_0\}.
  \]

  \item Fit linear models separately:
  \begin{itemize}[leftmargin=*]
    \item Left:
    \[
      y_i \approx m_1 x_i + b_1, \quad i \in \mathcal{L}
    \]
    using ordinary least squares (\texttt{np.polyfit}).

    \item Right:
    \[
      y_i \approx m_2 x_i + b_2, \quad i \in \mathcal{R}.
    \]
  \end{itemize}

  \item Construct predictions $y_i^{\mathrm{pred}}$ piecewise; compute
  \[
    \mathrm{SSE}(x_0) = \sum_i \left(y_i^{\mathrm{pred}} - y_i\right)^2.
  \]

  \item Choose $x_0$ that minimises SSE; record best $m_1, b_1, m_2, b_2$.
\end{enumerate}

So the final best parameters are
\[
  (m_1^\ast, b_1^\ast, m_2^\ast, b_2^\ast, x_0^\ast).
\]

\subsubsection*{Application to PSPEC: \texttt{joint\_piecewise\_linear\_fit}}

\begin{verbatim}
def joint_piecewise_linear_fit(uvp_power_list, uvpspec_avg_power_list,
                               method='trf', manual_fallback=True):
    xs, ys = [], []
    for pcoh, pinc in zip(uvp_power_list, uvpspec_avg_power_list):
        mask = np.isfinite(pcoh) & np.isfinite(pinc) & (pcoh>0) & (pinc>0)
        xs.append(np.log10(pcoh[mask]))
        ys.append(np.log10(pinc[mask]))
    x_all = np.concatenate(xs)
    y_all = np.concatenate(ys)

    return fit_piecewise_linear(x_all, y_all,
                                method=method,
                                manual_fallback=manual_fallback), x_all
\end{verbatim}

Here we set
\[
  x = \log_{10} P_{\mathrm{coh}}, \quad
  y = \log_{10} P_{\mathrm{inc}}.
\]

Model:
\[
  \log_{10} P_{\mathrm{inc}} \approx
  \begin{cases}
    m_1 \log_{10} P_{\mathrm{coh}} + b_1,
    & \log_{10} P_{\mathrm{coh}} \le x_0, \\[3pt]
    m_2 \log_{10} P_{\mathrm{coh}} + b_2,
    & \log_{10} P_{\mathrm{coh}} > x_0.
  \end{cases}
\]

\textbf{Interpretation:}

\begin{itemize}[leftmargin=*]
  \item At low coherent power (low SNR), the relationship is described by slope $m_1$.
  \item At high coherent power, slope $m_2$ tells how quickly $P_{\mathrm{inc}}$ grows with $P_{\mathrm{coh}}$.
\end{itemize}

If there were \emph{no} signal loss and no weird behaviour, you’d expect:
\[
  m_1 \approx 1, \quad m_2 \approx 1.
\]

A slope $m_2 < 1$ in the high-power regime means:
\[
  \log P_{\mathrm{inc}} \approx m_2 \log P_{\mathrm{coh}} + b_2
  \quad\Rightarrow\quad
  P_{\mathrm{inc}} \propto P_{\mathrm{coh}}^{\,m_2},
\]
which is ``sub-linear'' growth: large coherent power is \emph{compressed} when seen in the incoherent estimator -- i.e.\ a sign of multiplicative signal loss. The code turns this into a ``loss percent'' like $(1-m_2)\times 100\%$.

\section*{3. Generic linear combo model}

\subsection*{Function: \texttt{fit\_linear\_combo}}

This is a simple linear regression
\[
  y \approx a \cdot f(x) + b \cdot g(x) + c
\]
implemented as:

\begin{verbatim}
X = np.column_stack([fx, gx])  # (or [fx,gx,snx])
coefs = np.dot(np.linalg.pinv(X), y)
\end{verbatim}

In full form (when you include an intercept, or in your commented version):

\begin{itemize}[leftmargin=*]
  \item Data:
  \[
    y_i, f_i, g_i.
  \]

  \item Design matrix:
  \[
    X =
    \begin{bmatrix}
      f_1 & g_1 & 1 \\
      f_2 & g_2 & 1 \\
      \vdots & \vdots & \vdots \\
      f_N & g_N & 1 \\
    \end{bmatrix}.
  \]

  \item Coefficient vector:
  \[
    \boldsymbol{\beta} =
    \begin{bmatrix}
      a \\ b \\ c
    \end{bmatrix}.
  \]

  \item Model:
  \[
    \mathbf{y} \approx X \boldsymbol{\beta}.
  \]

  \item Least-squares solution (using pseudoinverse):
  \[
    \hat{\boldsymbol{\beta}} = X^+ \mathbf{y}
    \approx (X^\top X)^{-1} X^\top \mathbf{y}.
  \]
\end{itemize}

You use this style in the ``single fit'' toy block (currently disabled) to model
\[
  P_{\mathrm{inc}} \approx N_a\, P_N + m\, P_{\mathrm{coh}}.
\]

\section*{4. Joint PN model: globally fitting \texorpdfstring{$P_{\mathrm{inc}}$}{P\_inc} vs \texorpdfstring{$P_{\mathrm{coh}}$}{P\_coh} and \texorpdfstring{$P_N$}{P\_N} across all groups}

\subsection*{Function: \texttt{fit\_joint\_PN\_linear\_combo}}

This is your main ``PN model'' for signal loss.

\subsubsection*{Step 1 -- build arrays per group}

\begin{verbatim}
Ys, PN_cols, coh_cols, mask = [], [], [], []
for up, ua, un in zip(uvp_power_list,
                      uvpspec_avg_power_list,
                      uvpspec_noise_list):
    m = np.isfinite(ua)&np.isfinite(un)&np.isfinite(up)
    Ys.append(ua[m])   # P_inc
    PN_cols.append(un[m])   # P_N
    coh_cols.append(up[m])  # P_coh
    mask.append(m)
\end{verbatim}

For each group $g$:

\[
  Y^{(g)}_i = P_{\mathrm{inc}}^{(g)}(t_i), \quad
  PN^{(g)}_i = P_{N}^{(g)}(t_i), \quad
  C^{(g)}_i = P_{\mathrm{coh}}^{(g)}(t_i)
\]

with only finite points kept.

Then concatenate over groups:

\begin{verbatim}
Y = np.concatenate(Ys)           # shape (N,)
C = np.concatenate(coh_cols)     # shape (N,)
Ntot = Y.size                    # total number of data points
\end{verbatim}

So $Y$ is the stacked $P_{\mathrm{inc}}$ for all groups.

\subsubsection*{Step 2 -- ``block-diagonal'' PN columns}

\begin{verbatim}
X_cols = []
start = 0
for PN in PN_cols:
    Ni = PN.size
    col_full = np.zeros(Ntot, dtype=float)
    col_full[start:start+Ni] = PN
    X_cols.append(col_full)
    start += Ni
\end{verbatim}

This is key. For each group $g$:

\begin{itemize}[leftmargin=*]
  \item Let its length be $N_g$.
  \item You create a column vector of length $N$ filled with zeros except over the rows corresponding to that group, where you insert that group’s noise vector $PN^{(g)}$.
\end{itemize}

Schematically for two groups:
\[
  X_{\mathrm{PN}} =
  \begin{bmatrix}
    PN^{(1)}_1 & 0 \\
    PN^{(1)}_2 & 0 \\
    \vdots      & \vdots \\
    PN^{(1)}_{N_1} & 0 \\
    0 & PN^{(2)}_1 \\
    0 & PN^{(2)}_2 \\
    \vdots & \vdots \\
    0 & PN^{(2)}_{N_2}
  \end{bmatrix}.
\]

So group 1’s noise lives in column 1, group 2’s noise in column 2, etc. That’s why it behaves like a block-diagonal structure.

\subsubsection*{Step 3 -- build full design matrix}

\begin{verbatim}
X = np.column_stack(X_cols + [C])
\end{verbatim}

So the design matrix $X$ has:

\begin{itemize}[leftmargin=*]
  \item First $G$ columns = block PN columns,
  \item Last column = concatenated coherent power $C$.
\end{itemize}

Dimension:

\[
  X \in \mathbb{R}^{N \times (G+1)}.
\]

\subsubsection*{Step 4 -- fit coefficients}

\begin{verbatim}
coefs    = np.linalg.pinv(X) @ Y
Na_list  = coefs[:-1]
slopeloss = coefs[-1]
Y_pred = X @ coefs
\end{verbatim}

Let
\[
  \boldsymbol{\beta} =
  \begin{bmatrix}
    N_a^{(1)} \\ N_a^{(2)} \\ \vdots \\ N_a^{(G)} \\ m
  \end{bmatrix}.
\]

Then the model is:
\[
  Y \approx X \boldsymbol{\beta}.
\]

Unstacking by group, for each group $g$ and time $t_i$:
\[
  P_{\mathrm{inc}}^{(g)}(t_i) \approx
  N_a^{(g)}\, P_{N}^{(g)}(t_i) \;+\; m\, P_{\mathrm{coh}}^{(g)}(t_i).
\]

So mathematically:

\begin{itemize}[leftmargin=*]
  \item $N_a^{(g)}$ is a \emph{per-group scalar amplitude} for the noise power.
  \item $m = \text{slopeloss}$ is a \emph{global scalar} relating coherent and incoherent power after subtracting the noise contribution.
\end{itemize}

You solve this via least squares:
\[
  \hat{\boldsymbol{\beta}} = X^+ Y.
\]

\textbf{Physical interpretation:}

\begin{itemize}[leftmargin=*]
  \item If the pipeline were perfect and we had properly subtracted noise, you’d expect something like
  \[
    P_{\mathrm{inc}} \approx P_{\mathrm{coh}}
    \quad \Rightarrow \quad m \approx 1
  \]
  (and $N_a^{(g)}$ capturing the actual noise normalisation for each group).

  \item If the pipeline introduces multiplicative signal loss, the \emph{coherent} power transferred into the incoherent estimate is suppressed. The coefficient $m$ encodes that:
  \begin{itemize}[leftmargin=*]
    \item If $m > 1$, then
      \[
        P_{\mathrm{inc}} \approx N_a P_N + m P_{\mathrm{coh}},
      \]
      and you interpret $1/m$ as a ``loss ratio'' $P_{\text{true}} / P_{\text{measured}}$.
    \item In your code, the ``loss percentage'' is computed as:
      \[
        \text{loss} \approx (1 - 1/m)\times 100\%.
      \]
  \end{itemize}
\end{itemize}

\textbf{Use in plotting:}

\begin{verbatim}
median_vals.append((1-1/slopeloss_joint)*100)
...
ax_group[0].axhline(1/slopeloss_joint, label="All Group PN Fit Loss Ratio = 1/slopeloss")
\end{verbatim}

You then plot e.g.\ $P_{\mathrm{coh}} / P_{\mathrm{inc,fit}}$ vs.\ SNR, and compare that to the global constant $1/m$. If the ratio clusters around $1/m$, that’s your effective signal-loss factor.

\section*{5. PN + PSN model: adding a ``signal$\times$noise'' term per group}

\subsection*{Function: \texttt{fit\_joint\_PSN\_linear\_combo}}

Same structure but with extra columns for $P_{SN}^{(g)}$.

\subsubsection*{Step 1 -- build lists}

\begin{verbatim}
Ys, PN_cols, PSN_cols, coh_cols = [], [], [], []
for up, ua, un, usn in zip(uvp_power_list,
                           uvpspec_avg_power_list,
                           uvpspec_noise_list,
                           uvpspec_signalnoise_list):
    m = np.isfinite(ua)&np.isfinite(un)&np.isfinite(up)
    Ys.append(ua[m])       # P_inc
    PN_cols.append(un[m])  # P_N
    PSN_cols.append(usn[m])# P_SN (signal-noise combo)
    coh_cols.append(up[m]) # P_coh
\end{verbatim}

\subsubsection*{Step 2 -- concatenate and build block-diagonal columns}

You do the same block-diagonal trick for both \texttt{PN\_cols} and \texttt{PSN\_cols}:

\begin{verbatim}
X_cols  = block_diag_for(PN_cols)
XS_cols = block_diag_for(PSN_cols)
X = np.column_stack(X_cols + XS_cols + [C])
\end{verbatim}

So now:

\begin{itemize}[leftmargin=*]
  \item For each group you have two per-group scalars: one multiplying $P_N$, one multiplying $P_{SN}$.
  \item And one global slope $m$ multiplying $P_{\mathrm{coh}}$.
\end{itemize}

Let the parameter vector be:
\[
  \boldsymbol{\beta} =
  \begin{bmatrix}
    N_a^{(1)} \\ \dots \\ N_a^{(G)} \\
    N_b^{(1)} \\ \dots \\ N_b^{(G)} \\
    m
  \end{bmatrix}.
\]

For each group $g$ and time $t_i$:
\[
  P_{\mathrm{inc}}^{(g)}(t_i) \approx
  N_a^{(g)} P_N^{(g)}(t_i)
  + N_b^{(g)} P_{SN}^{(g)}(t_i)
  + m P_{\mathrm{coh}}^{(g)}(t_i).
\]

\textbf{Physical meaning of $P_{SN}$:}

You define:
\begin{verbatim}
uvpspec_signalnoise_list.append(
    np.sqrt( np.sqrt(2)*ups*un + (un)**2 )
)
\end{verbatim}
with $ups = P_S$, $un = P_N$. Inside the sqrt:
\[
  \sqrt{2}\, P_S P_N + P_N^2
\]
This looks like a term from the variance of a cross-power estimator where noise mixes with signal (signal $\times$ noise terms) plus pure noise variance. So $P_{SN}$ is effectively a \emph{signal-noise coupling scale}.

So the PSN model allows each group its own weighting for ``pure noise'' and ``signal$\times$noise'' contributions, and still solves for a single global signal-loss slope $m$.

\section*{6. SNR model: \texorpdfstring{$P_{\mathrm{inc}}$}{P\_inc} as function of SNR and \texorpdfstring{$P_{\mathrm{coh}}$}{P\_coh}}

\subsection*{Function: \texttt{fit\_joint\_SNR\_linear\_combo}}

Same pattern, but now the per-group regressor is SNR instead of PN:

\begin{verbatim}
Ys, SNR_cols, coh_cols = [], [], []
for up, ua, usn in zip(uvp_power_list,
                       uvpspec_avg_power_list,
                       uvpspec_SNR_list):
    m = np.isfinite(ua)&np.isfinite(usn)&np.isfinite(up)
    Ys.append(ua[m])        # P_inc
    SNR_cols.append(usn[m]) # SNR(t)
    coh_cols.append(up[m])  # P_coh
\end{verbatim}

Block-diagonal:

\begin{verbatim}
X_cols = block_diag_for(SNR_cols)
X = np.column_stack(X_cols + [C])
coefs = np.linalg.pinv(X) @ Y
SNRa_list = coefs[:-1]
slopeloss = coefs[-1]
\end{verbatim}

Model per group:
\[
  P_{\mathrm{inc}}^{(g)}(t_i) \approx
  \alpha^{(g)} \, \mathrm{SNR}^{(g)}(t_i)
  + m \, P_{\mathrm{coh}}^{(g)}(t_i).
\]

Here $\alpha^{(g)}$ is per-group scaling of SNR, and $m$ again is the global signal-loss slope.

This model tests whether \emph{explicitly} including SNR as a regressor helps explain the mapping from coherent to incoherent PSPEC beyond just a linear relationship in $P_{\mathrm{coh}}$ and $P_N$.

(In your current config \texttt{model\_SNR = False}, so it’s not used in the plots, but the math is there.)

\section*{7. No-PN model: direct slope between \texorpdfstring{$P_{\mathrm{inc}}$}{P\_inc} and \texorpdfstring{$P_{\mathrm{coh}}$}{P\_coh} only}

\subsection*{Function: \texttt{fit\_joint\_noPN\_linear\_combo}}

This is the simplest global model:

\begin{verbatim}
Ys, coh_cols = [], []
for up, ua in zip(uvp_power_list, uvpspec_avg_power_list):
    m = np.isfinite(ua)&np.isfinite(up)
    Ys.append(ua[m])
    coh_cols.append(up[m])
Y = np.concatenate(Ys)
C = np.concatenate(coh_cols)

X = np.column_stack([C])
coefs = np.linalg.pinv(X) @ Y
slopeloss = coefs[-1]  # a scalar
Y_pred = X @ coefs
\end{verbatim}

Here:
\[
  P_{\mathrm{inc}}(t) \approx m \, P_{\mathrm{coh}}(t),
\]
with one global scalar $m$ and no noise terms.

The least-squares solution is the classic slope:
\[
  m = \frac{\sum_t P_{\mathrm{coh}}(t) P_{\mathrm{inc}}(t)}
           {\sum_t P_{\mathrm{coh}}(t)^2}.
\]

\textbf{Physical meaning:} this lumps noise into residuals. It gives the simplest ``compression factor'' between coherent and incoherent power across all groups. Comparing $m$ from this model to the PN model’s $m$ tells you how much explicitly modelling the noise changes your inferred loss.

\section*{8. Goodness-of-fit \& parameter covariance}

\subsection*{Function: \texttt{goodness\_of\_fit(X, Y, Y\_pred)}}

Given design matrix $X$, data $Y$, and prediction $Y_{\text{pred}}$:

\begin{itemize}[leftmargin=*]
  \item Residuals:
  \[
    r_i = Y_i - Y_{\text{pred},i}.
  \]

  \item Sum of squared errors (SSE):
  \[
    \mathrm{SSE} = \sum_i r_i^2.
  \]

  \item Total sum of squares:
  \[
    \mathrm{SST} = \sum_i (Y_i - \bar{Y})^2,
    \quad \bar{Y} = \frac{1}{N}\sum_i Y_i.
  \]

  \item Mean squared error (MSE):
  \[
    \mathrm{MSE} = \frac{\mathrm{SSE}}{N}.
  \]

  \item Root MSE:
  \[
    \mathrm{RMSE} = \sqrt{\mathrm{MSE}}.
  \]

  \item Mean absolute error:
  \[
    \mathrm{MAE} = \frac{1}{N}\sum_i |r_i|.
  \]

  \item Coefficient of determination:
  \[
    R^2 = 1 - \frac{\mathrm{SSE}}{\mathrm{SST}}.
  \]
\end{itemize}

For parameter covariance, with $X \in \mathbb{R}^{N\times p}$:

\begin{verbatim}
MSE_samp = SSE / (N - p)
XtX = X.T @ X
inv_XtX = np.linalg.pinv(XtX)
cov_beta = MSE_samp * inv_XtX
\end{verbatim}

This is the usual linear-regression result:
\[
  \widehat{\mathrm{Cov}}(\hat{\boldsymbol{\beta}}) = \hat{\sigma}^2 (X^\top X)^{-1},
  \quad
  \hat{\sigma}^2 = \frac{\mathrm{SSE}}{N-p}.
\]

You then take:

\begin{itemize}[leftmargin=*]
  \item \texttt{PN\_cov\_beta[-1, -1]} $\rightarrow$ variance of the \emph{slope-loss} coefficient $m$ in the PN model.
  \item $\sigma_m = \sqrt{\mathrm{Var}(m)}$ is its 1-$\sigma$ uncertainty.
\end{itemize}

This lets you quote error bars on the loss ratio $1/m$ and the corresponding percentage loss.

\section*{9. Putting it together: how these models diagnose PSPEC signal loss}

Within \texttt{plot\_percent\_difference\_with\_1hrbin}, the key logic is:

\begin{enumerate}[leftmargin=*]
  \item \textbf{Collect data across all baselines for a given SPW \& pol}

\begin{verbatim}
*_, lst_r, up, ua, un, ups, usnr, _ = get_plot_data_all_blpairs(...)
\end{verbatim}

So you get:

\begin{itemize}[leftmargin=*]
  \item $P_{\mathrm{coh}}(t)$ (\texttt{up}),
  \item $P_{\mathrm{inc}}(t)$ (\texttt{ua}),
  \item $P_{N}(t)$ (\texttt{un}),
  \item $P_S(t)$ real part (\texttt{ups}),
  \item $\mathrm{SNR}(t)$ (\texttt{usnr}).
\end{itemize}

  \item \textbf{Segment in LST if desired}

Using \texttt{bump\_n\_val\_segments\_by\_dips} on $P_{\mathrm{coh}}(t)$ and $P_{\mathrm{inc}}(t)$, you can identify LST ``bumps'' and valleys, which correspond to different sky regions or beam footprints. That allows you to see if loss is LST-dependent.

  \item \textbf{Two-line log--log model}

Using \texttt{joint\_piecewise\_linear\_fit} on $\log_{10} P_{\mathrm{coh}}$ vs.\ $\log_{10} P_{\mathrm{inc}}$, you estimate:

\begin{itemize}[leftmargin=*]
  \item A low-power slope $m_1$.
  \item A high-power slope $m_2$.
  \item Break $x_0$ corresponding to a particular coherent PS level.
\end{itemize}

The high-power slope $m_2 < 1$ signals compression of bright modes.

  \item \textbf{Joint PN model}

Using \texttt{fit\_joint\_PN\_linear\_combo}, you model:
\[
  P_{\mathrm{inc}}^{(g)}(t) \approx N_a^{(g)} P_N^{(g)}(t) + m P_{\mathrm{coh}}^{(g)}(t),
\]
with:

\begin{itemize}[leftmargin=*]
  \item Per-group $N_a^{(g)}$ to absorb noise normalisation.
  \item Global $m$ capturing signal-loss.
\end{itemize}

This is a physically motivated decomposition:

\begin{itemize}[leftmargin=*]
  \item $P_N^{(g)}$ encodes variance of thermal noise for that group and time.
  \item $P_{\mathrm{coh}}^{(g)}$ is your ``true'' sky signal as measured by the coherent estimator.
  \item Any mismatch between $P_{\mathrm{inc}}$ and the naive sum $P_{\mathrm{coh}} + P_N$ is absorbed into $m$.
\end{itemize}

If the pipeline attenuates the coherent modes before they show up in the incoherent average, the fitted $m$ will depart from 1, giving a quantitative measure of signal loss.

  \item \textbf{Alternative models (PSN, SNR)}

\begin{itemize}[leftmargin=*]
  \item \texttt{fit\_joint\_PSN\_linear\_combo} adds an extra per-group term $P_{SN}$ to test whether \emph{signal--noise coupling} improves the description of $P_{\mathrm{inc}}$.
  \item \texttt{fit\_joint\_SNR\_linear\_combo} uses SNR as a regressor, probing whether the mapping depends explicitly on SNR beyond what $P_{\mathrm{coh}}$ and $P_N$ provide.
\end{itemize}

These are experiments to see if more refined noise modelling changes the inferred loss slope $m$.

  \item \textbf{Error bars and diagnostics}

The \texttt{goodness\_of\_fit} function gives you:

\begin{itemize}[leftmargin=*]
  \item Overall quality of the fit (R$^2$, RMSE, MSE\_samp).
  \item Covariance of fitted parameters, especially var$(m)$.
\end{itemize}

That lets you say things like ``the loss factor $1/m$ is $0.85 \pm 0.03$'' and compare different models or subsets (different SPWs, LST ranges, or baseline groups).
\end{enumerate}

If you’d like, next step we can:

\begin{itemize}[leftmargin=*]
  \item Pick \emph{one} of these models (say the PN model) and write out the design matrix and solution explicitly for a toy example with 2 groups and 3 time samples each, just to make the block-diagonal structure and the $P_{\mathrm{inc}} \approx N_a P_N + m P_{\mathrm{coh}}$ relation absolutely concrete.
\end{itemize}

\end{document}
